\begin{figure}[htbp]
\centering
\begin{tikzpicture}[font=\small,>=Stealth]
  \fill[astroPale,rounded corners=4pt] (.2,.2) rectangle (6.05,4.1);
  \fill[astroPale,rounded corners=4pt] (6.35,.2) rectangle (12.2,4.1);
  \node[font=\bfseries,text=astroNavy] at (3.12,3.65){$h\nu<\Phi$};
  \node[font=\bfseries,text=astroNavy] at (9.27,3.65){$h\nu>\Phi$};
  \fill[astroNavy] (.55,.75) rectangle (5.7,1.15);
  \fill[astroNavy] (6.7,.75) rectangle (11.85,1.15);
  \foreach \x in {.9,1.7,2.5,3.3,4.1,4.9,5.35}{\fill[white] (\x,.95) circle (.055);}
  \foreach \x in {7.05,7.85,8.65,9.45,10.25,11.05,11.5}{\fill[white] (\x,.95) circle (.055);}
  \fill[astroGold] (3.3,1.3) circle (.11);
  \draw[astroCoral,line width=1.6pt] (3.07,1.55)--(3.53,2.01);
  \draw[astroCoral,line width=1.6pt] (3.53,1.55)--(3.07,2.01);
  \draw[->,astroGold,line width=1.7pt] (2.1,3.2)--(3.05,2.08);
  \node[astroInk!70] at (3.12,.45){electrón retenido};
  \draw[astroBlue,line width=1.8pt] (8.1,3.05) .. controls (8.35,3.35) and (8.6,2.75) .. (8.85,3.05) .. controls (9.1,3.35) and (9.35,2.75) .. (9.6,3.05);
  \draw[->,astroBlue,line width=1.7pt] (9.6,3.05)--(9.85,2.05);
  \fill[astroCoral] (9.87,2.0) circle (.11);
  \draw[->,astroCoral,line width=1.8pt] (9.87,2.0)--(10.65,2.85);
  \node[astroInk!70] at (9.25,.45){fotoelectrón emitido};
\end{tikzpicture}
\caption{Por debajo del umbral el electrón permanece ligado; por encima, puede escapar de la superficie.}
\end{figure}